\begin{figure}[t]
\centering
\begin{tikzpicture}[
  font=\small,
  cell/.style={draw=hzslate!45, rounded corners=2pt, minimum width=35mm,
               minimum height=14mm, align=center, inner sep=2pt},
  base/.style={cell, fill=hzslate!10},
  new/.style={cell, fill=hzblue!9, draw=hzblue!50},
  hd/.style={hzslate!85, font=\small\bfseries, align=center}
]
\newcommand{\gloss}[1]{{\scriptsize\color{hzslate!85}#1}}

\node[hd] at (2.05, 1.35) {unrestricted};
\node[hd] at (6.15, 1.35) {scoped to\\the subject};

\node[hd, anchor=east] at (0.05, 0.0)  {whole\\turn};
\node[hd, anchor=east] at (0.05,-1.75) {claim\\span};

\node[base] at (2.05, 0.0)  {\texttt{vector}\\[2pt]\gloss{the usual baseline}};
\node[new]  at (6.15, 0.0)  {\texttt{scoped}\\[2pt]\gloss{control}};
\node[new]  at (2.05,-1.75) {\texttt{span}\\[2pt]\gloss{control}};
\node[base] at (6.15,-1.75) {\texttt{graph}\\[2pt]\gloss{the system under test}};

\draw[hzslate!25, line width=0.3pt] (0.25,-0.875) -- (7.95,-0.875);
\draw[hzslate!25, line width=0.3pt] (4.10, 0.85)  -- (4.10,-2.60);
\end{tikzpicture}
\caption{The four indexes as a $2\times2$: the unit that is indexed down the
side, what a query may reach across the top. \arm{vector} and \arm{graph} are
the comparison usually drawn, and they differ in both factors at once. The two
middle cells are the controls that say which factor a difference belongs to.
Every cell ranks the same weighted terms by cosine through the same store, is
read by the same reader, and is marked by the same scorer, so a difference
between cells is a difference of index and of nothing else.}
\label{fig:design}
\end{figure}
